
\section{Introducción}
El análisis de circuitos eléctricos complejos requiere de herramientas con capacidad de simplificación debido a la dificultad inherente presente en ellos. Un teorema es una proposición que parte de ciertas suposiciones o hipótesis para comprobar una tesis no evidente, siendo en este caso los teoremas de redes, los cuales resultan de gran utilidad y aplicación en el estudio, análisis y síntesis de circuitos. Emplear estas herramientas permite obtener soluciones simples y completas en una gran cantidad de redes que se encuentran en la práctica.

La aplicación de estos teoremas se fundamenta en la suposición subyacente de la unicidad de la solución de las redes eléctricas para un conjunto dado de entradas y condiciones iniciales. La razón primordial de su uso es que una red eléctrica modificada resulta significativamente más sencilla de estudiar y analizar que la red original. No obstante, es imperativo considerar que la generalidad y simplicidad de estos métodos pueden ser engañosas si se malinterpreta su significado o no se percibe correctamente el alcance de su aplicación.